\chapter{Collaborations et réponses aux besoins}
\label{ch:collaborations}

\questionchapitre{À quels besoins, exprimés par qui, ces développements répondent-ils
effectivement ?}

\section{Le cadre de travail}

Les développements présentés dans ce rapport ont été conçus et réalisés dans leur intégralité au
cours du postdoctorat, dans un dialogue continu avec trois interlocuteurs aux apports distincts.

\paragraph{Supervision scientifique : Nicolas Labroche (LIFAT).} Il a encadré l'orientation
scientifique des travaux, en particulier l'arbitrage
initial : accepter que le sujet d'explicabilité impose d'abord la construction de son terrain
expérimental, plutôt que de le traiter sur un jeu de données de circonstance. Cet arbitrage
détermine tout ce qui précède.

\paragraph{Collaboration technique : Hugo Breuillard (BRGM).} Data scientist en charge du projet
côté BRGM, il constitue le point de contact sur les questions de chaîne de données et de
stratégie de prévision. Les échanges ont porté sur l'articulation entre l'entrepôt développé ici
et les besoins de modélisation du programme, et sur la cohérence des indicateurs produits avec
ceux employés par ailleurs.

\paragraph{Expertise métier : Augustin Thomas (BRGM).} Hydrogéologue, il a été l'interlocuteur
sur l'ensemble des questions de domaine. C'est de ces échanges que procèdent les décisions qui
ancrent la plateforme dans la pratique hydrogéologique plutôt que dans une transposition
générique de méthodes d'apprentissage.

\section{La conduite du projet}

Une part du travail n'apparaît ni dans les dépôts ni dans les mesures, et il serait pourtant
trompeur de la passer sous silence : ces développements n'ont pas été commandés, ils ont été
proposés, puis portés jusqu'à leur adoption.

\paragraph{Identifier le besoin et le faire reconnaître.} L'absence de jeu de données
exploitable aurait pu être traitée comme une contrainte personnelle, contournée par une
extraction de circonstance suffisante pour produire un résultat. Le parti inverse a été pris :
formuler le manque, le documenter, et le soumettre aux partenaires pour vérifier s'il était
partagé. Il l'était. Cette vérification a transformé un obstacle individuel en objet de travail
légitime, et c'est elle qui a rendu possible l'investissement de plusieurs mois que représente
un entrepôt de données.

\paragraph{Aller chercher les interlocuteurs.} Les contacts n'ont pas été hérités d'un
dispositif préexistant. Ils ont été pris directement, auprès des interlocuteurs techniques comme
des experts métier, en préparant chaque échange autour des questions dont la réponse
conditionnait un choix de conception. La section suivante en donne l'exemple le plus net.

\paragraph{Tenir le suivi dans la durée.} Un échange ponctuel ne produit pas d'effet. Le suivi
avec le BRGM a été maintenu tout au long de la période, ce qui explique que les décisions métier
ne soient pas restées au stade de la discussion mais se retrouvent inscrites dans le code, comme
le montre la section~\ref{sec:apport-metier}.

\paragraph{Ouvrir le socle au-delà de son périmètre.} Enfin, les travaux ont été présentés aux
équipes du projet ANR \textsc{Aida}, dont le domaine d'expérimentation recouvre en partie celui
traité ici. L'intention n'était pas de revendiquer un périmètre, mais l'inverse : signaler
l'existence d'une infrastructure disponible, afin que ceux qui en auraient l'usage n'aient pas à
la reconstruire. Une infrastructure dont personne ne connaît l'existence ne vaut pas mieux
qu'une infrastructure inexistante.

\begin{center}
\begin{tabularx}{\textwidth}{@{}lX@{}}
\toprule
\textbf{Nature du travail} & \textbf{Ce qu'il a produit} \\
\midrule
Identification et formulation du besoin & Un constat vérifié auprès des partenaires, et non une
hypothèse personnelle \\
Prise de contact et préparation des échanges & Des décisions de conception fondées sur
l'expertise du domaine plutôt que sur des conventions d'informaticien \\
Suivi dans la durée & Des orientations métier effectivement traduites dans le code \\
Ouverture à d'autres équipes & Un socle connu au-delà de son projet d'origine \\
\bottomrule
\end{tabularx}
\end{center}

Cette dimension mérite d'être signalée pour une raison qui dépasse le compte rendu d'activité.
Un entrepôt de données mutualisé n'est pas un objet technique : c'est un objet d'organisation.
Sa valeur tient à ce que plusieurs équipes acceptent de s'y référer, ce qui relève de la
conviction et non de la programmation. Le chapitre~\ref{ch:socle} montre que les obstacles
restants à son élargissement sont eux aussi de cette nature.

\section{Ce que la collaboration métier a produit}
\label{sec:apport-metier}

L'apport d'expertise n'est pas resté au stade de la discussion : il est inscrit dans le code, et
identifiable. Quatre exemples le montrent.

\paragraph{Les configurations de modélisation par famille d'aquifère.} La table présentée au
chapitre~\ref{ch:metier} associe à chaque nature de milieu une combinaison de fonction de
réponse et de modèle de bruit. Cette connaissance ne s'obtient pas par validation croisée : elle
procède de la compréhension du comportement hydrodynamique de chaque milieu. Le cas karstique,
traité par une double exponentielle pour représenter les deux temps de réponse des conduits et
de la matrice, en est l'illustration la plus nette.

\paragraph{L'adoption d'un standard d'évaluation reconnu.} Le choix d'évaluer les modèles selon
les quatre critères STOWA plutôt que par un indicateur d'erreur unique relève du même registre.
Il introduit notamment le critère de temps de réponse rapporté à la période de calibration, qui
écarte des modèles apparemment excellents mais non identifiables, subtilité que la seule
pratique de l'apprentissage automatique ne fait pas apparaître.

\paragraph{La règle d'affichage cartographique.} La règle énoncée au
chapitre~\ref{ch:observatoire}, un indicateur réellement standardisé ou une valeur réelle et
jamais un intermédiaire, procède d'une préoccupation d'interprétation propre au domaine. Un
outil qui colorie une carte de France selon des seuils choisis par son auteur fabrique une norme
que l'utilisateur lira comme une référence établie.

\paragraph{L'arbitrage sur l'évapotranspiration.} La décision documentée en
section~\ref{sec:hargreaves} est exemplaire du bénéfice de la confrontation disciplinaire. La
variable fournie par la réanalyse était disponible, correctement nommée et directement
utilisable ; seule la confrontation de ses valeurs aux ordres de grandeur connus pour la France a
révélé qu'elle ne désignait pas la grandeur attendue. Aucune vérification interne au code ne
pouvait détecter cette erreur, car il ne s'agissait pas d'une erreur de calcul.

\section{Les cas d'usage identifiés avec le BRGM}

Les échanges n'ont pas seulement produit des décisions techniques. Ils ont fait émerger une
liste de besoins concrets, qui a servi à orienter les développements et qui reste disponible
pour la suite. Elle mérite d'être consignée, car elle constitue à elle seule un livrable : un
socle sans usage identifié n'est qu'une infrastructure en attente.

\begin{center}
\begin{tabularx}{\textwidth}{@{}lXl@{}}
\toprule
\textbf{Besoin exprimé} & \textbf{Ce qu'il recouvre} & \textbf{État} \\
\midrule
Qualifier les piézomètres & Distinguer automatiquement les ouvrages inertiels, réactifs et
influencés par des prélèvements, pour savoir lesquels sont exploitables & outillé \\
Prévoir les niveaux & Anticiper l'évolution d'une nappe à partir du climat & outillé \\
Expliquer les prévisions & Identifier les facteurs déterminants d'une prévision donnée &
outillé \\
Comparer les approches & Confronter méthodes statistiques, modèles métier et apprentissage
profond sur les mêmes données & possible, non fait \\
Explorer les données & Permettre aux experts de parcourir et de comparer les chroniques sans
écrire de requête & outillé \\
Scénariser la sécheresse & Répondre à des questions du type « quelle pluie faudrait-il pour
sortir de cet état ? », en appui des comités sécheresse & outillé, non éprouvé \\
Étudier les échanges nappe-rivière & Détecter et caractériser les relations entre une nappe et
le cours d'eau voisin, notamment le soutien d'étiage & non traité \\
Caractériser la double porosité & Distinguer, au sein d'une même formation, les écoulements
rapides des écoulements lents, que les méthodes actuelles peinent à séparer & non traité \\
Couvrir les nappes captives & Disposer d'un indicateur de situation pour les nappes qui n'en ont
pas, l'indicateur existant ne valant que pour les nappes libres & non traité \\
\bottomrule
\end{tabularx}
\end{center}

Deux remarques sur ce tableau. La colonne d'état est délibérément sévère : « outillé » signifie
que les moyens existent, pas qu'un résultat a été produit. Et les trois dernières lignes sont
assumées comme telles : ce sont des sujets à part entière, identifiés comme prometteurs au cours
des échanges, mais qui dépassaient le périmètre et le temps disponibles. Ils sont consignés ici
parce qu'ils constituent des points d'entrée pour la suite, et parce qu'ils viennent des experts
du domaine plutôt que d'une projection d'informaticien.

Un point mérite d'être souligné à propos de la scénarisation. Les équipes du BRGM travaillant
sur la situation des nappes emploient déjà les indicateurs standardisés que produit désormais
l'entrepôt. Coupler ces indicateurs à des scénarios prospectifs, c'est-à-dire répondre non
seulement « où en est-on » mais « que faudrait-il pour en sortir », est le prolongement naturel
de leur pratique actuelle. C'est cette continuité qui rend le socle mobilisable par eux sans
changement de méthode.

\section{L'articulation avec les travaux de recherche en cours}

Les développements présentés au chapitre~\ref{ch:xai} ne sont pas isolés. Des travaux menés au
LIFAT sous la direction de Nicolas Labroche portent sur les explications contrefactuelles et sur
la mesure de leur qualité en tant qu'explications, question ouverte que les métriques usuelles ne
tranchent pas.

La contribution de ce postdoctorat à cette thématique est \textbf{d'ordre technique}, et c'est
délibéré : fournir les données, les modèles, les méthodes comparables et les moyens de
validation sans lesquels cette recherche resterait cantonnée à des jeux d'essai génériques. La
répartition est claire, et il vaut mieux l'écrire que la laisser deviner.

\begin{center}
\begin{tabularx}{\textwidth}{@{}lX@{}}
\toprule
\textbf{Relève de ces travaux} & Le socle de données, les modèles à expliquer, l'implémentation
des méthodes, la référence de comparaison, le moyen de validation indépendant, et la
qualification des cas d'étude \\
\textbf{Relève des travaux en cours} & La définition de ce qu'est une bonne explication, les
métriques qui la mesurent, et l'évaluation comparative des approches \\
\bottomrule
\end{tabularx}
\end{center}

C'est aussi ce qui explique qu'aucune conclusion méthodologique ne figure dans ce rapport sur la
qualité des contrefactuels produits. Elle n'aurait pas eu de sens : la question relève d'un
travail en cours, et l'anticiper reviendrait à préjuger de son résultat.

\section{Des besoins distincts, un outil commun}

Les deux communautés visées n'attendent pas la même chose, et l'architecture a dû accommoder les
deux sans se scinder.

\begin{center}
\begin{tabularx}{\textwidth}{@{}p{3.6cm}XX@{}}
\toprule
 & \textbf{Hydrogéologues du BRGM} & \textbf{Chercheurs en apprentissage} \\
\midrule
\textbf{Question posée} & Où en est la ressource, et que se passe-t-il si l'on prélève ? &
Quelle méthode prédit le mieux, et pourquoi ? \\
\textbf{Attente principale} & Des indicateurs conformes aux pratiques, une scénarisation
crédible & Un jeu de données consolidé, un protocole reproductible \\
\textbf{Obstacle levé} & L'agrégation manuelle de sources hétérogènes & Le coût d'entrée à
l'expérimentation \\
\textbf{Ce qui a été livré} & Observatoire, indicateurs validés, laboratoire de modèles métier,
scénarios bornés & Entrepôt interrogeable, catalogue d'architectures, traçabilité des
expériences, outillage d'explicabilité \\
\bottomrule
\end{tabularx}
\end{center}

Le point de rencontre est l'entrepôt. Parce que les indicateurs métier et les jeux
d'apprentissage sont produits par la même chaîne et validés une seule fois, une comparaison
entre un modèle appris et un modèle métier porte réellement sur les modèles, et non sur des
préparations de données divergentes. Cette propriété, invisible dans l'interface, est ce qui rend
la plateforme utilisable comme instrument de recherche et non seulement comme outil de
consultation.

\section{État de mise à disposition}

La plateforme est déployée, l'interface sur l'infrastructure mutualisée de l'établissement, le
calcul et les bases sur un serveur équipé d'un processeur graphique. Son usage demeure à ce
jour circonscrit au cercle du projet. L'accès est nominatif : il n'existe ni inscription libre ni
authentification déléguée, les comptes étant créés par un administrateur. Ce choix, adapté à un
outil scientifique manipulant une chaîne de traitement en évolution, constitue le principal frein
à une diffusion plus large, et le premier levier à actionner pour l'obtenir.

La collaboration a donc effectivement transformé les développements, et non seulement accompagné
leur réalisation : les décisions métier structurantes sont inscrites dans le code et
documentées, et les deux communautés disposent d'un socle commun de données.

Plusieurs pistes identifiées au cours des échanges n'ont pas été explorées, notamment le
rapprochement avec les modèles à réservoirs du BRGM et l'accès à leurs travaux de comparaison de
modèles. Quant à l'élargissement des usages au-delà du cercle du projet, il relève d'une
politique d'accès et d'un accompagnement davantage que de développements supplémentaires.
